\documentclass[11pt,a4paper]{article}
\usepackage[utf8]{inputenc}
\usepackage[T1]{fontenc}
\usepackage[french]{babel}
\usepackage[margin=1.5cm]{geometry}
\usepackage{titlesec}
\usepackage{enumitem}
\usepackage{xcolor}
\usepackage{amsfonts}

% Configuration des titres : Ajout de "blue" souligné avant le nom du chapitre
\titleformat{\section}{\large\bfseries\sffamily}{}{0pt}{\color{blue}blue \underline}
\titleformat{\subsection}{\normalsize\bfseries}{}{0pt}{\textit}

\title{Fiche de Révision Complète - CEJM}
\author{Maël - BTS SIO}
\date{}

\begin{document}

\maketitle

\section{Les agents économiques et leurs relations}
\subsection{Les types d'entreprises par activité}
\begin{itemize}[label=$\bullet$]
    \item \textbf{Entreprise Industrielle} : Activité de transformation. Elle achète des matières premières pour fabriquer des produits finis (ex: usine, construction).
    \item \textbf{Entreprise Commerciale (Marchande)} : Activité de négoce. Elle achète des marchandises et les revend en l'état sans transformation (ex: Fnac, Carrefour).
    \item \textbf{Entreprise de Services} : Fournit une prestation immatérielle. On distingue les services marchands (payants) et non marchands (publics).
\end{itemize}

\subsection{Rappels sur les agents et les flux}
\begin{itemize}[label=$\bullet$]
    \item \textbf{Ménages} : Fonction principale de consommation.
    \item \textbf{Administrations publiques (APU)} : Produisent des services non marchands et redistribuent les revenus.
    \item \textbf{Flux économiques} : Le \textbf{flux réel} (transfert de bien/travail) a pour contrepartie un \textbf{flux monétaire} (paiement/salaire).
\end{itemize}

\section{Le fonctionnement des marchés}
\begin{itemize}[label=$\bullet$]
    \item \textbf{Loi de l'offre et de la demande} : Détermine le prix d'équilibre. Si l'Offre < Demande, le prix augmente (pénurie).
    \item \textbf{Asymétrie d'information} : Situation où l'un des partenaires dispose d'informations que l'autre n'a pas (risque de sélection adverse).
    \item \textbf{Externalités} : Impact d'un agent sur un autre sans transaction monétaire. L'État intervient par la fiscalité (écotaxes) ou des incitations (subventions).
\end{itemize}

\section{La formation du contrat et sa validité}
\subsection{Les vices du consentement}
Pour qu'un contrat soit valide, le consentement doit être libre et éclairé. S'il est vicié, le contrat peut être annulé :
\begin{itemize}[label=--]
    \item \textbf{L'Erreur} : Se tromper sur un élément essentiel du contrat (la qualité de l'objet ou la personne).
    \item \textbf{Le Dol} : Tromperie intentionnelle, manœuvres frauduleuses ou silence sur une info déterminante pour pousser l'autre à signer.
    \item \textbf{La Violence} : Pression physique ou morale (menace) exercée pour contraindre à la signature.
\end{itemize}

\subsection{Autres conditions de validité}
\begin{itemize}[label=--]
    \item \textbf{La Capacité juridique} : Aptitude à avoir des droits et à les exercer (être majeur et non protégé).
    \item \textbf{Le Contenu licite et certain} : L'objet du contrat doit exister et ne pas être contraire à l'ordre public.
\end{itemize}

\section{La création d'entreprise et les logiques}
\begin{itemize}[label=$\bullet$]
    \item \textbf{Logique Entrepreneuriale} : Capacité à identifier des opportunités, innover et prendre des risques (Schumpeter).
    \item \textbf{Logique Managériale} : Organisation, planification et optimisation de l'existant pour assurer la pérennité.
    \item \textbf{Complémentarité} : Le développement d'une structure nécessite souvent de passer de l'intuition de l'entrepreneur à la rigueur du manager.
\end{itemize}

\section{Les finalités de l'entreprise et parties prenantes}
\begin{itemize}[label=$\bullet$]
    \item \textbf{Finalités} : Économique (profit/croissance), Sociale (bien-être des salariés), Sociétale (RSE/impact environnemental).
    \item \textbf{Parties Prenantes} : Acteurs internes (salariés, actionnaires) ou externes (clients, État, ONG) qui influencent ou subissent les décisions de l'entreprise.
    \item \textbf{Gouvernance} : Arbitrage entre les intérêts parfois divergents des parties prenantes.
\end{itemize}

\section{La performance de l'entreprise}
La performance est multidimensionnelle : elle n'est plus seulement financière.
\begin{itemize}[label=$\bullet$]
    \item \textbf{Efficacité vs Efficience} : L'efficacité est l'atteinte des objectifs. L'efficience est l'atteinte des objectifs en optimisant les ressources.
    \item \textbf{Le Tableau de Bord Prospectif (TBP)} : Outil de pilotage stratégique via 4 axes :
    \begin{enumerate}
        \item \textit{Axe Financier} : Rentabilité, valeur actionnariale.
        \item \textit{Axe Client} : Satisfaction, fidélisation, part de marché.
        \item \textit{Axe Processus Internes} : Qualité, délais, innovation technique.
        \item \textit{Axe Apprentissage organisationnel} : Compétences, motivation des salariés, climat social.
    \end{enumerate}
\end{itemize}

\section{L'influence de l'État et la Régulation}
L'État intervient pour réguler l'économie et sécuriser les relations.
\begin{itemize}[label=$\bullet$]
    \item \textbf{Régulation du marché} : L'État lutte contre les abus de position dominante pour garantir une concurrence saine.
    \item \textbf{Politiques structurelles} : Actions de long terme (éducation, infrastructures, recherche).
    \item \textbf{Évolutions législatives} : La \textbf{Loi PACTE (2019)} redéfinit l'entreprise. Elle introduit la "Raison d'être" et permet de devenir une \textbf{Société à mission} (poursuivre des objectifs sociaux/environnementaux en plus du profit).
\end{itemize}

\section{L'implication des ressources humaines}
Les RH sont un facteur de compétitivité majeur.
\begin{itemize}[label=$\bullet$]
    \item \textbf{Gestion des compétences} : Adéquation entre les besoins de l'entreprise et les capacités des salariés (Savoir, Savoir-faire, Savoir-être).
    \item \textbf{Motivation et Implication} : Dépendent de facteurs matériels (salaire) et psychologiques (reconnaissance, autonomie).
    \item \textbf{Le Processus de Management (Fayol)} :
    \begin{itemize}
        \item \textit{Prévoir} : Fixer les objectifs et la stratégie.
        \item \textit{Organiser} : Allouer les ressources.
        \item \textit{Commander} : Piloter et donner les directives.
        \item \textit{Coordonner} : Créer de la synergie entre les services.
        \item \textit{Contrôler} : Vérifier les résultats et corriger les écarts.
    \end{itemize}
\end{itemize}

\end{document}